El CVRP ha sido objeto de estudio sistemático en la literatura de investigación
operacional desde su introducción a fines de los años cincuenta
\parencite{dantzig1959truck}. Su carácter $\mathcal{NP}$-hard
\parencite{lenstra1981complexity} explica que, junto a los métodos exactos revisados en
el Informe~1, se haya desarrollado un amplio cuerpo de \textit{métodos aproximados}
capaces de escalar a instancias industriales. En este informe el alcance se centra en
las metaheurísticas y, en particular, en \textit{Particle Swarm Optimization}.

\subsection{Metaheurísticas para el CVRP}

Las metaheurísticas suelen clasificarse en métodos basados en \textit{trayectoria}, que
operan sobre una única solución que se modifica iterativamente (como \textit{Tabu
Search} o \textit{Simulated Annealing}), y métodos basados en \textit{población}, que
mantienen y hacen evolucionar un conjunto de soluciones, como los algoritmos genéticos,
\textit{Ant Colony Optimization} (ACO) y \textit{Particle Swarm Optimization} (PSO).
La literatura de metaheurísticas para el CVRP es extensa y abarca desde algoritmos
constructivos con búsqueda local hasta esquemas poblacionales sofisticados. Entre los
métodos de trayectoria destacan la búsqueda tabú y el recocido simulado; entre los
poblacionales, los algoritmos genéticos, ACO y PSO; y en años recientes, los esquemas de
búsqueda de vecindario amplio (LNS/ALNS) figuran entre los enfoques aproximados más
competitivos para las variantes del VRP \parencite{tan2021vehicle}. Todos estos métodos
renuncian a la garantía de optimalidad a cambio de escalar a instancias de tamaño
industrial, y con frecuencia se combinan con búsqueda local para intensificar la
exploración del espacio de soluciones.

El presente trabajo se sitúa dentro de los métodos poblacionales, adoptando PSO. Una
revisión específica de su aplicación al problema de ruteo, con una comparación de sus
principales variantes, se encuentra en \textcite{marinakis2018particle}, que sirve de
referencia para los apartados siguientes.

\subsection{Particle Swarm Optimization}

PSO fue introducido por \textcite{kennedy1995particle} como un algoritmo de optimización
inspirado en el comportamiento social de bandadas de aves y cardúmenes de peces. Cada
\textit{partícula} representa una solución candidata que se desplaza por el espacio de
búsqueda con una \textit{velocidad} actualizada en función de su mejor posición
histórica (\textit{pbest}) y de la mejor posición del enjambre o del vecindario
(\textit{gbest}). Contribuciones posteriores refinaron su dinámica: el
\textit{peso de inercia} \parencite{shi1998modified} y el \textit{factor de constricción}
\parencite{clerc2002particle} regulan el equilibrio entre exploración y explotación y
garantizan estabilidad de la convergencia. Una síntesis unificada de estas variantes se
encuentra en \textcite{poli2007particle}.

% NOTA: La mecánica detallada de PSO (ecuaciones de velocidad y posición) va en el
% Capítulo 3 (Método), no aquí. Aquí basta el posicionamiento en la literatura.

\subsection{Adaptación de PSO a problemas combinatorios y al CVRP}
\label{subsec:adaptacion_pso}

La formulación canónica de PSO opera sobre un espacio continuo $\mathbb{R}^d$, mientras
que el CVRP es un problema combinatorio cuyas soluciones son permutaciones y particiones
de clientes. Como las nociones de velocidad y dirección carecen de un significado natural
sobre recorridos y permutaciones \parencite{poli2007particle}, aplicar PSO al CVRP exige
un mecanismo que relacione las posiciones continuas del enjambre con soluciones de ruteo
factibles. La literatura ha abordado esta brecha por tres vías principales.

\subsubsection{Codificación de valor real y decodificación}

La vía más difundida conserva la dinámica continua de PSO y delega la traducción al
espacio combinatorio a un procedimiento de \textit{decodificación}: la posición continua
de una partícula se interpreta, típicamente ordenando sus componentes, como una lista de
prioridad o una permutación de clientes, que luego se particiona en rutas respetando la
capacidad. \textcite{ai2009particle} propusieron dos representaciones de referencia para
el CVRP. La primera, SR-1, es un vector de $n+2m$ dimensiones cuyas primeras $n$
componentes inducen una lista de prioridad de clientes (por orden ascendente de su
valor) y cuyas últimas $2m$ componentes fijan puntos de referencia que orientan la
asignación de clientes a vehículos; la segunda, SR-2, extiende esta idea con una
representación alternativa. Ambas incorporan procedimientos de mejora local en la
decodificación para elevar la calidad de las soluciones. Sobre esta misma línea,
\textcite{ahmed2018bilayer} introdujeron un método de decodificación novedoso combinado
con una búsqueda local en dos capas (BLS-PSO), que alcanza las mejores soluciones
conocidas en la mayoría de las instancias de prueba; su resultado ilustra el papel
determinante que cumplen tanto la decodificación como la hibridación con búsqueda local
en el desempeño final.

\subsubsection{PSO discreto}

Una segunda vía redefine los operadores de PSO directamente sobre el espacio discreto,
reinterpretando la velocidad como una secuencia de intercambios (\textit{swaps}) que
transforma una permutación en otra. \textcite{chen2006hybrid} desarrollaron un PSO
discreto (DPSO) para el CVRP que combina búsqueda global y local y se hibrida con
recocido simulado para escapar de óptimos locales, reportando buen desempeño en
instancias de gran tamaño. En una línea afín, \textcite{mauliddina2020implementation}
implementaron un DPSO para el CVRP y ajustaron sus parámetros mediante un diseño
experimental de un factor a la vez (OFAT), contrastando el resultado con modelos de PSO
propuestos previamente.

\subsubsection{Enfoques híbridos}

Una tercera vía combina PSO con otras metaheurísticas o con esquemas de agrupamiento.
\textcite{kao2012hybrid} propusieron un híbrido de ACO y PSO para el CVRP en el que cada
hormiga memoriza su mejor solución, al estilo de una partícula. \textcite{son2021capacitated}
adoptaron un enfoque en dos fases (agrupamiento de clientes mediante $k$-Means seguido de
una fase de ruteo con un híbrido adaptativo de PSO y \textit{Grey Wolf Optimization}),
obteniendo soluciones cercanas a las de los métodos exactos. Estas variantes evidencian
una tendencia general: el PSO puro rara vez alcanza por sí solo las mejores soluciones
conocidas, por lo que la hibridación con búsqueda local u otros mecanismos de
intensificación es habitual en los mejores resultados reportados para el CVRP.

\subsection{Posicionamiento del enfoque adoptado}

El Informe~1 mostró que el esquema exacto certifica optimalidad sobre instancias pequeñas
y medianas, pero pierde tratabilidad al crecer el tamaño del problema: el salto de $n=32$
a $n=33$ bastó para desplazar la instancia fuera del rango resoluble a optimalidad dentro
del presupuesto de tiempo. Este límite motiva el uso de una metaheurística, capaz de
obtener soluciones de buena calidad en tiempos acotados sobre instancias que exceden el
alcance del método exacto. Se adopta PSO por tratarse de un método poblacional
ampliamente estudiado, de implementación relativamente simple y con una dinámica de
búsqueda bien caracterizada \parencite{poli2007particle}.

De las tres vías de adaptación revisadas en la Sección~\ref{subsec:adaptacion_pso}, este
trabajo adopta la codificación de valor real con decodificación, siguiendo la
representación SR de \textcite{ai2009particle}. La elección responde a tres razones.
Primero, preserva la dinámica continua de PSO, de modo que las ecuaciones canónicas de
velocidad y posición se aplican sin modificación, a diferencia del PSO discreto
\parencite{chen2006hybrid}, que exige redefinir la velocidad como secuencias de
intercambios. Segundo, constituye una representación de referencia para el CVRP, con un
procedimiento de decodificación documentado que garantiza la factibilidad de capacidad de
forma constructiva. Tercero, evita la complejidad de los esquemas híbridos o en dos fases
\parencite{son2021capacitated,kao2012hybrid}, cuya infraestructura excede el alcance de
este trabajo y dificultaría aislar el comportamiento propio de PSO.

La revisión también muestra que los mejores resultados, como la obtención de las mejores
soluciones conocidas, suelen provenir de variantes hibridadas con búsqueda local
\parencite{ahmed2018bilayer}. No obstante, el objetivo de este informe no es competir por
la mejor solución conocida, sino caracterizar el compromiso entre el método exacto y una
metaheurística representativa; por ello se opta por una implementación de PSO con
representación SR, transparente y bien comprendida, cuyo detalle se desarrolla en el
Capítulo~\ref{cap:metodo}. La incorporación de búsqueda local se contempla como una
posible extensión.